```latex
\documentclass[11pt]{article}

\usepackage[a4paper,margin=1in]{geometry}
\usepackage{amsmath,amsthm,amssymb,mathtools}

\newtheorem{theorem}{Theorem}
\newtheorem{proposition}[theorem]{Proposition}
\newtheorem{corollary}[theorem]{Corollary}
\newtheorem{remark}[theorem]{Remark}
\newtheorem{definition}[theorem]{Definition}

\newcommand{\kk}{\Bbbk}
\newcommand{\Uc}{\mathcal U}
\newcommand{\bcS}{\mathcal S}
\newcommand{\fb}{\mathfrak b}
\newcommand{\fn}{\mathfrak n}
\newcommand{\fB}{\mathbb B}
\newcommand{\bipr}{\mathbin{\bowtie}}

\begin{document}

\title{The punctured bigon, the generalized quantum double, and \(\Uc_q(\mathfrak{sl}_n)\)}
\author{}
\date{}
\maketitle

\paragraph{Conventions.}
Fix \(n\ge 2\), and work over the field
\[
\kk=\mathbb C(q^{1/n}).
\]
Let \(\Uc_q(\mathfrak{sl}_n)\) denote the simply connected form of the Drinfeld--Jimbo quantum group (the algebra denoted \(U_q^{\mathrm{ext}}(\mathfrak{sl}_n)\) in Klimyk--Schm\"udgen), with generators
\[
E_i,\ F_i\quad (1\le i\le n-1),\qquad \hat K_j^{\pm1}\quad (1\le j\le n),
\]
subject to the usual relations and \(\hat K_1\cdots \hat K_n=1\).  As usual,
\[
K_i:=\hat K_i\hat K_{i+1}^{-1}\qquad (1\le i\le n-1).
\]
Let \(\Uc_q(\fb_+)\) and \(\Uc_q(\fb_-)\) be the positive and negative Borel Hopf subalgebras.

We write \(e_{ij}\) for the standard matrix units in \(\operatorname{End}(\kk^n)\), and we fix the standard type-\(A_{n-1}\) \(R\)-matrix
\begin{equation}\label{eq:R}
R
=
q\sum_{i=1}^n e_{ii}\otimes e_{ii}
+
\sum_{i\neq j} e_{ii}\otimes e_{jj}
+
(q-q^{-1})\sum_{i<j} e_{ij}\otimes e_{ji}.
\end{equation}
Any nonzero scalar multiple of \(R\) yields the same RTT/RLL relations, so the scalar
\(q^{-1/n}\) attached to the defining \(\Uc_q(\mathfrak{sl}_n)\)-module is immaterial here.

For \(X\in \operatorname{End}(\kk^n)\otimes A\), we write
\[
X_1:=X\otimes 1,\qquad X_2:=1\otimes X
\]
inside \(\operatorname{End}(\kk^n\otimes \kk^n)\otimes A\).

\medskip

The proof is a comparison of two \emph{standard} presentations of the same triangular
\(RLL\)-algebra:
\begin{itemize}
\item the punctured-bigon presentation of \(\bcS_n(\fB^\ast)\) by the two horizontal families of arcs;
\item the Faddeev--Reshetikhin--Takhtajan presentation of the generalized quantum double of the Borels.
\end{itemize}

\begin{definition}\label{def:An}
Let \(\mathcal A_n\) be the \(\kk\)-algebra generated by symbols
\[
b_{ij}\qquad (1\le i\le j\le n),\qquad
c_{ij}\qquad (1\le j\le i\le n),
\]
with each \(b_{ii}\) and \(c_{ii}\) invertible, assembled into an upper triangular matrix
\[
B:=\sum_{i\le j} e_{ij}\otimes b_{ij}
\]
and a lower triangular matrix
\[
C:=\sum_{i\ge j} e_{ij}\otimes c_{ij},
\]
subject to
\begin{equation}\label{eq:RLL}
RB_1B_2=B_2B_1R,\qquad
RC_1C_2=C_2C_1R,\qquad
RB_1C_2=C_2B_1R,
\end{equation}
and the special-linear relations
\begin{equation}\label{eq:det1}
b_{11}b_{22}\cdots b_{nn}=1,\qquad
c_{11}c_{22}\cdots c_{nn}=1.
\end{equation}
\end{definition}

\begin{remark}
The choice of convention in \eqref{eq:RLL} is the one for which the upper horizontal family
in the punctured bigon is identified with the positive \(L\)-operator and the lower horizontal family with the negative \(L\)-operator.
Equivalent conventions obtained by replacing \(R\) by \(R_{21}^{-1}\) and simultaneously transposing the matrix indexing lead to the same final result.
\end{remark}

\begin{proposition}\label{prop:two-presentations}
The algebra \(\mathcal A_n\) admits the following two realizations.

\smallskip

\noindent\textup{(a)} The punctured-bigon stated skein algebra \(\bcS_n(\fB^\ast)\) is canonically isomorphic to \(\mathcal A_n\), via the identification of the generators \(b_{ij},c_{ij}\) with the two right-going horizontal arcs (above and below the puncture).

\smallskip

\noindent\textup{(b)} Let
\[
\mathfrak D:=\Uc_q(\fb_-)\bipr \Uc_q(\fb_+)
\]
be the generalized quantum double defined by the Drinfeld--Tanisaki skew pairing
\[
\langle \hat K_i,\hat K_j\rangle=q^{\frac1n-\delta_{ij}},
\qquad
\langle E_i,F_j\rangle=\frac{\delta_{ij}}{q^{-1}-q},
\qquad
\langle \hat K_i,F_j\rangle=\langle E_j,\hat K_i\rangle=0.
\]
If \(L^+=(\ell^+_{ij})\) is the standard upper triangular FRT matrix and
\(L^-=(\ell^-_{ij})\) the standard lower triangular FRT matrix in \(\mathfrak D\),
then the assignment
\[
b_{ij}\longmapsto \ell^+_{ij},\qquad
c_{ij}\longmapsto \ell^-_{ij}
\]
extends to an algebra isomorphism \(\mathcal A_n\xrightarrow{\sim}\mathfrak D\).
Moreover, if \(\hat K_i^+\) and \(\hat K_i^-\) denote the Cartan generators in the positive and negative Borel copies inside \(\mathfrak D\), then
\begin{equation}\label{eq:simple-coeffs}
\ell^+_{ii}=(\hat K_i^+)^{-1},\qquad
\ell^-_{ii}=\hat K_i^-,
\end{equation}
and
\begin{equation}\label{eq:simple-offdiag}
\ell^+_{i,i+1}=(q^{-1}-q)(\hat K_i^+)^{-1}E_i,\qquad
\ell^-_{i+1,i}=(q-q^{-1})F_i\hat K_i^-.
\end{equation}
\end{proposition}

\begin{proof}
Part (a) is exactly the standard punctured-bigon presentation theorem for the higher stated skein algebra:
the three matrix equations in \eqref{eq:RLL} are the three local crossing resolutions for two right-going horizontal arcs, and \eqref{eq:det1} is the skein-theoretic expression of the fact that the \(n\)-labeled strand is the tensor unit in the \(\mathfrak{sl}_n\)-spider.  We use this presentation theorem as the skein-theoretic input.

For part (b), this is the type-\(A\) FRT realization of the generalized quantum double of the two Borels; see Faddeev--Reshetikhin--Takhtajan and Klimyk--Schm\"udgen, Chapter~8.  Concretely, the generalized double \(\mathfrak D\) contains a pair of triangular matrices \(L^\pm\) whose coefficients satisfy \eqref{eq:RLL}, and because we are in the simply connected \(\mathfrak{sl}_n\)-form, their diagonal coefficients satisfy \eqref{eq:det1}.  The displayed formulas \eqref{eq:simple-coeffs} and \eqref{eq:simple-offdiag} are the standard formulas for the diagonal and nearest off-diagonal matrix coefficients in that realization.  In particular, these coefficients generate the two Borel factors, hence generate \(\mathfrak D\).  The FRT theorem asserts that there are no further algebraic relations beyond \eqref{eq:RLL} and \eqref{eq:det1}.  Therefore \(b_{ij}\mapsto \ell^+_{ij}\), \(c_{ij}\mapsto \ell^-_{ij}\) yields an isomorphism \(\mathcal A_n\cong \mathfrak D\).
\end{proof}

\begin{theorem}\label{thm:main}
There is a unique algebra isomorphism
\[
\Phi:\bcS_n(\fB^\ast)\xrightarrow{\sim}\Uc_q(\fb_-)\bipr \Uc_q(\fb_+)
\]
such that, if \(B=(b_{ij})\) and \(C=(c_{ij})\) are the two horizontal-arc matrices in \(\bcS_n(\fB^\ast)\), then
\[
\Phi(B)=L^+,\qquad \Phi(C)=L^-.
\]
Equivalently,
\[
\Phi(b_{ij})=\ell^+_{ij},\qquad \Phi(c_{ij})=\ell^-_{ij}.
\]

In particular,
\[
\Phi(b_{ii})=(\hat K_i^+)^{-1},\qquad
\Phi(c_{ii})=\hat K_i^-,
\]
and
\[
\Phi(b_{i,i+1})=(q^{-1}-q)(\hat K_i^+)^{-1}E_i,\qquad
\Phi(c_{i+1,i})=(q-q^{-1})F_i\hat K_i^-.
\]
\end{theorem}

\begin{proof}
Compose the two isomorphisms in Proposition~\ref{prop:two-presentations}:
\[
\bcS_n(\fB^\ast)\xrightarrow{\ \sim\ }\mathcal A_n\xrightarrow{\ \sim\ }\Uc_q(\fb_-)\bipr \Uc_q(\fb_+).
\]
This gives the required algebra isomorphism, and the displayed formulas are exactly
\eqref{eq:simple-coeffs}--\eqref{eq:simple-offdiag}.
\end{proof}

\begin{remark}\label{rem:inverse}
Theorem~\ref{thm:main} implies, in particular, that the diagonal punctured-bigon generators are invertible. Solving the formulas in Theorem~\ref{thm:main} for the standard generators of the two Borels gives
\[
(\hat K_i^+)^{-1}\longleftrightarrow b_{ii},\qquad
\hat K_i^-\longleftrightarrow c_{ii},
\]
\[
E_i\longleftrightarrow \frac{1}{q^{-1}-q}\,b_{ii}^{-1}b_{i,i+1},
\qquad
F_i\longleftrightarrow \frac{1}{q-q^{-1}}\,c_{i+1,i}c_{ii}^{-1}.
\]
Thus the punctured-bigon generators recover the standard Chevalley generators explicitly.
\end{remark}

\begin{corollary}\label{cor:quotient}
Let \(I\) be the two-sided ideal of \(\bcS_n(\fB^\ast)\) generated by
\[
b_{ii}c_{ii}-1\qquad (1\le i\le n).
\]
Then there is an algebra isomorphism
\[
\bcS_n(\fB^\ast)/I \xrightarrow{\sim} \Uc_q(\mathfrak{sl}_n).
\]

More explicitly, if \(\bar b_{ij},\bar c_{ij}\) denote the images of \(b_{ij},c_{ij}\) in the quotient, then under this isomorphism
\[
\bar b_{ii}\longmapsto \hat K_i^{-1},\qquad
\bar c_{ii}\longmapsto \hat K_i,
\]
\[
\bar b_{i,i+1}\longmapsto (q^{-1}-q)\hat K_i^{-1}E_i,\qquad
\bar c_{i+1,i}\longmapsto (q-q^{-1})F_i\hat K_i.
\]
Equivalently,
\[
\hat K_i \longleftrightarrow \bar c_{ii}=\bar b_{ii}^{-1},
\qquad
E_i \longleftrightarrow \frac{1}{q^{-1}-q}\,\bar c_{ii}\bar b_{i,i+1},
\qquad
F_i \longleftrightarrow \frac{1}{q-q^{-1}}\,\bar c_{i+1,i}\bar b_{ii}.
\]
\end{corollary}

\begin{proof}
Let
\[
\mathfrak D:=\Uc_q(\fb_-)\bipr \Uc_q(\fb_+),
\]
and let \(J\subset \mathfrak D\) be the two-sided ideal generated by
\[
\ell^+_{ii}\ell^-_{ii}-1\qquad (1\le i\le n).
\]
By Theorem~\ref{thm:main},
\[
\Phi(b_{ii}c_{ii}-1)=\ell^+_{ii}\ell^-_{ii}-1,
\]
hence \(\Phi(I)=J\), and therefore
\[
\bcS_n(\fB^\ast)/I \xrightarrow{\sim} \mathfrak D/J.
\]

It remains to show that \(\mathfrak D/J\cong \Uc_q(\mathfrak{sl}_n)\).

\smallskip

Write \(\hat K_i^\pm\) for the Cartan generators in the two Borel copies inside \(\mathfrak D\). By \eqref{eq:simple-coeffs},
\[
\ell^+_{ii}\ell^-_{ii}-1
=
(\hat K_i^+)^{-1}\hat K_i^- -1.
\]
Hence \(J\) is precisely the ideal identifying the two Cartan copies:
in the quotient \(\mathfrak D/J\), the images of \(\hat K_i^+\) and \(\hat K_i^-\) become equal; denote their common image by \(\hat K_i\).

Now use the PBW/triangular decomposition.  Let
\[
T^\pm:=\kk[(\hat K_1^\pm)^{\pm1},\dots,(\hat K_n^\pm)^{\pm1}]
\big/\big((\hat K_1^\pm\cdots \hat K_n^\pm)-1\big),
\]
and
\[
T:=\kk[\hat K_1^{\pm1},\dots,\hat K_n^{\pm1}]
\big/\big((\hat K_1\cdots \hat K_n)-1\big).
\]
By the standard PBW theorem for quantum Borels,
\[
\Uc_q(\fb_-)\cong \Uc_q(\fn_-)\otimes T^-,
\qquad
\Uc_q(\fb_+)\cong T^+\otimes \Uc_q(\fn_+)
\]
as vector spaces. Since \(\mathfrak D\) is, by definition of the generalized double, the vector space tensor product
\[
\mathfrak D \cong \Uc_q(\fb_-)\otimes \Uc_q(\fb_+),
\]
we obtain a vector space factorization
\[
\mathfrak D \cong \Uc_q(\fn_-)\otimes T^-\otimes T^+\otimes \Uc_q(\fn_+).
\]
Because the generators of \(J\) lie entirely in the middle torus factor and these torus elements normalize the nilpotent parts, quotienting by \(J\) gives
\[
\mathfrak D/J \cong \Uc_q(\fn_-)\otimes T\otimes \Uc_q(\fn_+)
\]
as vector spaces.

On the other hand, the standard triangular decomposition of the simply connected quantum group gives
\[
\Uc_q(\mathfrak{sl}_n)\cong \Uc_q(\fn_-)\otimes T\otimes \Uc_q(\fn_+)
\]
as vector spaces as well.

Finally, multiplication of the two Borel subalgebras inside \(\Uc_q(\mathfrak{sl}_n)\) and the universal property of the generalized double give a surjective algebra homomorphism
\[
\pi:\mathfrak D\longrightarrow \Uc_q(\mathfrak{sl}_n)
\]
sending the positive and negative Borel copies to the corresponding subalgebras of \(\Uc_q(\mathfrak{sl}_n)\). Since the two Cartan copies already coincide in \(\Uc_q(\mathfrak{sl}_n)\), \(\pi\) kills \(J\), hence factors through a surjective algebra map
\[
\bar\pi:\mathfrak D/J\longrightarrow \Uc_q(\mathfrak{sl}_n).
\]
The two preceding vector-space decompositions show that \(\bar\pi\) is between two vector spaces with the same PBW basis parametrization; therefore \(\bar\pi\) is bijective. Thus
\[
\mathfrak D/J\cong \Uc_q(\mathfrak{sl}_n).
\]

Composing with \(\bcS_n(\fB^\ast)/I\cong \mathfrak D/J\) gives the stated isomorphism. The explicit formulas are immediate from Theorem~\ref{thm:main} after identifying \(\hat K_i^+=\hat K_i^-=\hat K_i\) in the quotient.
\end{proof}

\begin{remark}
If one prefers the usual simple-root toral generators \(K_r=\hat K_r\hat K_{r+1}^{-1}\) (\(1\le r\le n-1\)), then Corollary~\ref{cor:quotient} reads
\[
K_r\longleftrightarrow
\bar c_{rr}\,\bar c_{r+1,r+1}^{-1}
=
\bar b_{r+1,r+1}\,\bar b_{rr}^{-1},
\]
\[
E_r\longleftrightarrow \frac{1}{q^{-1}-q}\,\bar c_{rr}\bar b_{r,r+1},
\qquad
F_r\longleftrightarrow \frac{1}{q-q^{-1}}\,\bar c_{r+1,r}\bar b_{rr}.
\]
Thus the quotient of the punctured-bigon algebra recovers the standard Drinfeld--Jimbo generators explicitly.
\end{remark}

\begin{thebibliography}{9}

\bibitem{FRT}
L.~D.~Faddeev, N.~Yu.~Reshetikhin, and L.~A.~Takhtajan,
\emph{Quantization of Lie groups and Lie algebras},
Leningrad Math. J. \textbf{1} (1990), 193--225.

\bibitem{KS}
A.~Klimyk and K.~Schm\"udgen,
\emph{Quantum Groups and Their Representations},
Texts and Monographs in Physics, Springer, Berlin, 1997.

\end{thebibliography}

\end{document}
```